% ============================================================
% 20_tables.tex — Tabellen-Helfer (semantisches Table-System)
% ============================================================
% Zentrales Design-Prinzip:
%   1. Tabellen fuellen IMMER \linewidth (kein manuelles p{0.16..}).
%   2. Spaltenbreiten werden PROPORTIONAL deklariert (Y{0.18} L{0.54} ...),
%      sodass die Summe der Faktoren = Anzahl X-Spalten ergibt.
%   3. Zebra-Streifen + Header-Highlight sind ins Environment eingebaut;
%      direkte Farb-Primitive sind in Kapiteln per Linter verboten.
%
% Damit ist strukturell ausgeschlossen:
%   - zu kurze Zebra-Streifen (Tabelle schmaler als Box),
%   - überstehende Zeilen (Content sprengt Box-Breite),
%   - Drift zwischen Kapiteln (jede Tabelle nutzt dieselben Wrapper).

% ------------------------------------------------------------
% Spaltentypen (tabularx-X-Varianten)
% ------------------------------------------------------------
% Feste Varianten — jede zaehlt als Faktor 1:
% m{#1}: X-Spalten vertikal mittig (Text + Grafik in einer Zeile).
\renewcommand{\tabularxcolumn}[1]{m{#1}}
\newcolumntype{L}{>{\RaggedRight\arraybackslash}X}   % linksbuendig, umbricht
\newcolumntype{C}{>{\Centering\arraybackslash}X}     % zentriert,  umbricht
\newcolumntype{R}{>{\RaggedLeft\arraybackslash}X}    % rechtsbuendig, umbricht

% Parametrisierte, proportionale Varianten — #1 = relativer Anteil.
%
% WICHTIG (tabularx-Konvention): Die Summe aller Faktoren in EINER
% Tabelle muss = Anzahl X-Spalten sein, damit die Tabelle \linewidth
% voll ausfuellt. L/C/R zählen als Faktor 1.
%
% Beispiele (2 Spalten → Summe 2):
%   Y{0.6}  Y{1.4}            → 30% / 70%
%   Y{1}    Y{1}              → 50% / 50% (== L L)
%   Q{0.8}  Y{1.2}            → 40% / 60%, Spalte 1 rechtsbuendig
% Beispiele (3 Spalten → Summe 3):
%   Y{0.6}  Y{1.5}  Y{0.9}    → 20% / 50% / 30%
%   L       Y{2}    L         → 25% / 50% / 25%
% Beispiele (4 Spalten → Summe 4):
%   Y{1.2}  Y{0.8}  Y{1.2}  Y{0.8}  → 30% / 20% / 30% / 20%
\newcolumntype{Y}[1]{>{\hsize=#1\hsize\RaggedRight\arraybackslash}X} % linksb., proport.
\newcolumntype{Z}[1]{>{\hsize=#1\hsize\Centering\arraybackslash}X}   % zentriert, proport.
\newcolumntype{Q}[1]{>{\hsize=#1\hsize\RaggedLeft\arraybackslash}X}  % rechtsb., proport.

% Zentrierte Formel-Spalte (Math-Mode automatisch, proportional)
\newcolumntype{F}[1]{>{\hsize=#1\hsize\Centering\arraybackslash$}X<{$}}

% ------------------------------------------------------------
% Luft zwischen Zellinhalt und Zeilenkante
% ------------------------------------------------------------
% Der Regler `rows` versprach »mehr Höhe für Zellen mit Brüchen« und konnte
% das nicht halten: \arraystretch skaliert das Zeilenstreben, und ein Streben
% ist ein MINDESTmass. Inhalt, der höher ist als es — ein \dfrac, eine Wurzel,
% ein Bild —, bleibt davon unberührt. Gemessen: an einer Textabelle +22 %
% Höhe, an derselben Tabelle mit \dfrac-Zellen +5 %.
%
% Bündig steht der Inhalt zusätzlich, und zwar unten. Grund ist die
% Ausrichtung der m-Spalte in array.sty (\ar@align@mcell): Ist eine Zelle
% höher als ein Streben, senkt sie den Block um
%   0.5 * (Zellhöhe - Strebenhöhe + \baselineskip)
% und schiebt ihn damit so weit nach unten, dass die Tiefe der Zelle die Tiefe
% der Zeile bestimmt — die ganze verbleibende Luft landet oben. Die Formel ist
% für mehrzeiligen Text gedacht; für eine einzelne hohe Zeile bedeutet sie
% Unterkante bündig. Im Satz sichtbar: Der Nenner eines \dfrac berührte die
% Zeilengrenze, in der letzten Zeile lief er aus der Box.
%
% Kein strebenbasierter Weg behebt das, weil jede Strebe ein Mindestmass ist
% und die Zeilentiefe das Maximum aus Strebe und Zellinhalt. Deshalb meldet
% der Ausrichter die Zelle um \ZSFtableCellPadY höher UND tiefer, ohne sie zu
% verschieben — dieselbe \raisebox-Geste wie bei \ZSFimage, nur eine Ebene
% höher und damit für JEDEN hohen Zellinhalt.
%
% Nur der hohe Zweig wird angefasst. Kurzer Text läuft weiter über die
% Grundlinien-Ausrichtung von array.sty; an ihm ändert sich nichts.
%
% \ar@align@mcell ist ein Hook und keine Interna-Bastelei: array.sty prüft ihn
% selbst mit \ifx\ar@align@mcell\@undefined und setzt ihn in anderen Pfaden
% auf \relax. Die Wache darunter hält den Fall offen, dass eine ältere
% array.sty ihn gar nicht kennt — dann bleibt alles beim Alten.
%
% Bewusst mit TeX-Primitiven und einem EIGENEN Register: \raisebox und die
% \@tempdim*-Register sind hier belegt — colortbl haelt den Farbueberhang der
% Zeile in \@tempdimb/\@tempdimc, und eine Zuweisung darauf wuerde die
% Zeilenflaeche verschieben. Mit \raisebox gemessen: die Zellen kamen leer
% heraus.
\newlength{\ZSFtableCellPadY}
\makeatletter
\newdimen\ZSF@cellShift
\@ifundefined{ar@align@mcell}{}{%
  \def\ar@align@mcell{%
    \ifdim \ht\ar@mcellbox > \ht\strutbox
      \ZSF@cellShift\ht\ar@mcellbox
      \advance\ZSF@cellShift-\ht\@arstrutbox
      \advance\ZSF@cellShift\baselineskip
      \ZSF@cellShift=.5\ZSF@cellShift
      \setbox\ar@mcellbox\hbox{\lower\ZSF@cellShift\box\ar@mcellbox}%
      \ht\ar@mcellbox\dimexpr\ht\ar@mcellbox+\ZSFtableCellPadY\relax
      \dp\ar@mcellbox\dimexpr\dp\ar@mcellbox+\ZSFtableCellPadY\relax
      \box\ar@mcellbox
    \else
      \box\ar@mcellbox
    \fi}%
}
\makeatother

% ------------------------------------------------------------
% Zebra-Streifen (Scannbarkeit)
% ------------------------------------------------------------
\newcommand{\ZSFzebraBG}{\TableZebraColor}
\newcommand{\SetZSFzebraBG}[1]{\renewcommand{\ZSFzebraBG}{#1}}
%
%
%
\newcommand{\ZSFzebraStart}[1]{\rowcolors{#1}{\BoxPlainBackColor}{\ZSFzebraBG}}
\newcommand{\ZSFzebraOff}{\rowcolors{1}{}{}}

% ------------------------------------------------------------
% Header-Highlight
% ------------------------------------------------------------
% Konvention:
%   \ZSFheaderRow   — startet die Header-Zeile und setzt den Zeilen-Hintergrund.
%   \ZSFhead{<Text>} — umschliesst den Inhalt JEDER Header-Zelle (Farbe + bold).
%
% Warum zweigeteilt? In tabularx-X-Spalten propagiert \color, das direkt nach
% \rowcolor am Zeilenanfang gesetzt wird, NICHT zuverlässig über die
% &-Grenzen. Die erste Zelle verliert die Farbe, andere behalten sie — je
% nach Spaltenbreite und Content (siehe SI-Praefix-Bug). Per-Zelle-Wrapper
% via \ZSFhead{...} ist robust und konsistent.
\newcommand{\ZSFheaderRowBG}{\chaptercolor}
\newcommand{\ZSFheaderRowFG}{\chaptercolortext}
\newcommand{\ZSFheaderRow}{\rowcolor{\ZSFheaderRowBG}}
% \ZSFInkOwned steht direkt neben dem \color, das den Kontrast setzt: Die
% Kopfzelle ist eine gesaettigte Flaeche mit eigener Textfarbe, eine
% Inline-Farbe darin waere unlesbar (40_colors_structure -> Ink-Vertrag).
\newcommand{\ZSFhead}[1]{{\color{\ZSFheaderRowFG}\ZSFInkOwned\ZSFfontTableHead #1}}

%
%
%
%
%
%
%
%
%
%
%
%
%
%
%
%
%
%
%
%
%
%
%
%
%
%
%
%
%
%
%
%
%
%
%
%
%
%
%
%
%
%
%
%
\makeatletter
\newcommand{\ZSF@spanBadAlign}[1]{%
  \PackageError{ZSF}{Unbekannte Ausrichtung '#1' in \string\ZSFspan}%
    {Gueltig sind L, C und R -- dieselben Buchstaben wie bei den
     Spaltentypen L/C/R.}%
}
\newcommand{\ZSFspan}[3]{%
  \if L#1\multicolumn{#2}{l}{#3}%
  \else\if R#1\multicolumn{#2}{r}{#3}%
  \else\if C#1\multicolumn{#2}{c}{#3}%
  \else\multicolumn{#2}{c}{\ZSF@spanBadAlign{#1}#3}%
  \fi\fi\fi
}
\makeatother

%
%
%
%
%
%
%
%
%
%
%
%
%
%
%
%
%
%
%
%
%
%
%
%
%
%
%
%
%
%
%
%
%
%
%
%
%
%
\makeatletter
\newif\ifZSF@tblHeader
\newif\ifZSF@tblZebra
\newcommand{\ZSF@tblFont}{\ZSFfontContent}
\newcommand{\ZSF@tblStretch}{\ZSFtableArrayStretch}
\newlength{\ZSF@tblColSep}

%
%
%
%
%
%
%
%
%
%
%
%
%
%
%
%
%
%
%
\newcommand{\ZSF@setRows}[1]{%
  \IfStrEqCase{#1}{%
    {normal}{\renewcommand{\ZSF@tblStretch}{\ZSFtableArrayStretch}%
             \setlength{\ZSFtableCellPadY}{\ZSFDensityTables\ZSFspaceXS}}%
    {roomy} {\renewcommand{\ZSF@tblStretch}{\ZSFtableArrayStretchRoomy}%
             \setlength{\ZSFtableCellPadY}{\ZSFDensityTables\ZSFspaceS}}%
    {tight} {\renewcommand{\ZSF@tblStretch}{\ZSFtableArrayStretchTight}%
             \setlength{\ZSFtableCellPadY}{0pt}}%
  }%
}
\newcommand{\ZSF@setColsep}[1]{%
  \IfStrEqCase{#1}{%
    {normal}{\setlength{\ZSF@tblColSep}{\ZSFtableEdgePad}}%
    {tight} {\setlength{\ZSF@tblColSep}{\ZSFtableEdgePadTight}}%
  }%
}
\newcommand{\ZSF@tblSpec}{}
\newcommand{\ZSF@beginTabularx}[1]{\tabularx{\linewidth}{#1}}
\pgfkeys{
  /zsf/table/.is family, /zsf/table,
  header/.is if=ZSF@tblHeader,
  zebra/.is if=ZSF@tblZebra,
  font/.is choice,
  font/normal/.code={\renewcommand{\ZSF@tblFont}{\ZSFfontContent}},
  font/dense/.code={\renewcommand{\ZSF@tblFont}{\ZSFfontContentDense}},
  % `rows` und `colsep` beantworten ihre Frage EINMAL — die beiden Schlüssel
  % hier reichen sie nur durch (\ZSF@setRows / \ZSF@setColsep, oben). Das
  % valuegrid stellt dieselben zwei Fragen und liest dieselbe Antwort; stünde
  % die Zuordnung an beiden Bausteinen, hiesse `rows=roomy` beim nächsten
  % Eingriff auf einem der beiden etwas anderes.
  rows/.is choice,
  rows/normal/.code={\ZSF@setRows{normal}},
  rows/roomy/.code ={\ZSF@setRows{roomy}},
  rows/tight/.code ={\ZSF@setRows{tight}},
  colsep/.is choice,
  colsep/normal/.code={\ZSF@setColsep{normal}},
  colsep/tight/.code ={\ZSF@setColsep{tight}},
  ZSF@tblDefaults/.style={header=true, zebra=true, font=normal,
                          rows=normal, colsep=normal},
}
\NewDocumentEnvironment{ZSFtable}{O{} m}{%
  \begingroup
  % Zellinhalt: ein \ZSFimage hier bekommt Zellen-, nicht Abbildungshoehe.
  % Explizit und nicht als Default verlassen, damit eine Tabelle IN einer
  % figbox nicht deren Abbildungsbudget erbt (siehe \ZSF@imgBudget, 60_boxes).
  \renewcommand{\ZSF@imgBudget}{\ZSFimageMaxHeight}%
  \pgfkeys{/zsf/table, ZSF@tblDefaults, #1}%
  \setlength{\tabcolsep}{\ZSF@tblColSep}%
  \renewcommand{\arraystretch}{\ZSF@tblStretch}%
  \ZSF@tblFont
  \ifZSF@tblZebra
    \ifZSF@tblHeader\ZSFzebraStart{2}\else\ZSFzebraStart{1}\fi
  \else
    \ZSFzebraOff
  \fi
%
%
%
%
%
%
%
%
%
%
%
%
%
%
%
%
%
%
%
%
%
%
%
%
%
%
%
%
%
%
%
%
%
%
%
  \edef\ZSF@tblSpec{@{\hskip\ZSF@tblColSep}#2@{\hskip\ZSF@tblColSep}}%
  \expandafter\ZSF@beginTabularx\expandafter{\ZSF@tblSpec}%
}{%
  \endtabularx
  \endgroup
}
\makeatother
